% from ega3-6-fr.tex
\oldpage[III]{164}
l'isomorphisme canonique
\[
  (M\otimes_B B')_{g'}
  \xrightarrow{\sim}
  M_g\otimes_{B_g}B'_{g'}
\]
($0_{\mathrm I}$, 1.5.4); nous laissons les détails au lecteur.
\end{env}

\begin{env}[6.7.11]
La théorie des hypertor globaux, développée ci-dessus pour deux
$S$-morphismes $X^{(i)}\to Y^{(i)}$ et deux complexes
$\mathscr P_\bullet^{(i)}$ de Modules quasi-cohérents limités inférieurement,
s'étend aussitôt au cas général suivant : on a un préschéma $S$, une famille
finie de $S$-préschémas $Y^{(i)}$ ($i\in I$), une famille finie de
$S$-morphismes séparés et quasi-compacts
$f_i:X^{(i)}\to Y^{(i)}$, et pour chaque $i$ un complexe de
$\mathscr O_{X^{(i)}}$-Modules quasi-cohérents
$\mathscr P_\bullet^{(i)}$ limités inférieurement. Si $Y$ est le produit des
$S$-préschémas $Y^{(i)}$, on définit alors pour chaque entier
$n\in\mathbf Z$, un $\mathscr O_Y$-Module quasi-cohérent
$\mathfrak{Tor}_n^S((f_i)_{i\in I};
(\mathscr P_\bullet^{(i)})_{i\in I})$,
ces Modules formant un $\partial$-foncteur covariant en chacun des complexes
$\mathscr P_\bullet^{(i)}$; en outre, ce foncteur est l'aboutissement commun
de six foncteurs spectraux
${}^{(t)}\mathscr E((f_i)_{i\in I};(\mathscr P_\bullet^{(i)})_{i\in I})$.
Nous laissons au lecteur le soin de répéter pour ce cas général les
définitions et les raisonnements faits ci-dessus pour $I=\{1,2\}$. Notons
simplement que lorsque $I$ se réduit à un seul élément, on retrouve
l'hypercohomologie $\mathscr H^\bullet(f,\mathscr P_\bullet)$ définie dans
(6.2.7) (comme on l'a déjà observé dans (6.7.4)). Lorsque $I$ est l'intervalle
$1\leq i\leq m$ de $\mathbf N$, nous écrirons
\[
  \mathfrak{Tor}_n^S
  (f_1,\ldots,f_m;
   \mathscr P_\bullet^{(1)},\ldots,\mathscr P_\bullet^{(m)})
  \quad\text{pour}\quad
  \mathfrak{Tor}_n^S
  ((f_i)_{i\in I};(\mathscr P_\bullet^{(i)})_{i\in I}).
\]
\end{env}

\begin{proposition}[6.7.12]
Les notations étant celles de (6.7.11), soit $J$ une partie de $I$ telle que,
pour $i\in I-J$, on ait $X^{(i)}=Y^{(i)}=S$, $f_i$ étant réduit à l'identité,
et $\mathscr P_\bullet^{(i)}$ égal au complexe réduit au terme de degré $0$
égal à $\mathscr O_S$. Il y a alors un isomorphisme canonique de
$\partial$-foncteurs
\[
  \mathfrak{Tor}_\bullet^S
  ((f_i)_{i\in I};(\mathscr P_\bullet^{(i)})_{i\in I})
  \xrightarrow{\sim}
  \mathfrak{Tor}_\bullet^S
  ((f_i)_{i\in J};(\mathscr P_\bullet^{(i)})_{i\in J}).
  \tag{6.7.12.1}\label{III.6.7.12.1.fr}
\]
\end{proposition}

\begin{proof}
On peut se borner à définir cet isomorphisme lorsque $S$ et les $Y^{(i)}$
sont affines, le recollement se faisant comme d'ordinaire. Pour $i\in I-J$,
on peut prendre le recouvrement $\mathfrak U^{(i)}$ formé du seul ensemble
$S$, et alors
$L_{\bullet\bullet}^{(i)}=C^\bullet(\mathfrak U^{(i)},
\mathscr P_\bullet^{(i)})$ est réduit à son seul terme de degrés $(0,0)$,
égal à $\Gamma(S,\mathscr O_S)=A(S)$; l'isomorphisme (6.7.12.1) est alors
évident.
\end{proof}

\begin{remark}[6.7.13]
Les notations étant celles de (6.7.3), considérons le $S$-isomorphisme
canonique
$Y^{(1)}\times_S Y^{(2)}\to Y^{(2)}\times_S Y^{(1)}$ (I, 3.3.5); alors
l'image par cet isomorphisme de
$\mathfrak{Tor}_\bullet^S(f_1,f_2;
\mathscr P_\bullet^{(1)},\mathscr P_\bullet^{(2)})$ est
$\mathfrak{Tor}_\bullet^S(f_2,f_1;
\mathscr P_\bullet^{(2)},\mathscr P_\bullet^{(1)})$; la question étant
locale, on est ramené au cas envisagé dans (6.6.2), et si on désigne par
$M_{\bullet\bullet\bullet}^{(i)}$ une résolution projective de
Cartan-Eilenberg de $L_{\bullet\bullet}^{(i)}$ ($i=1,2$), l'isomorphisme
considéré transforme
$M_{\bullet\bullet\bullet}^{(1)}\otimes_A
M_{\bullet\bullet\bullet}^{(2)}$ en
$M_{\bullet\bullet\bullet}^{(2)}\otimes_A
M_{\bullet\bullet\bullet}^{(1)}$, d'où notre assertion en considérant
l'homologie des complexes simples associés à ces tricomplexes.
\end{remark}

\subsection{Les suites spectrales d'associativité des hypertor globaux.}

\begin{env}[6.8.1]
Les hypothèses et notations étant celles de (6.7.11) (et en particulier les
$\mathscr P_\bullet^{(i)}$ étant supposés limités inférieurement, supposons
donnée une partition $(I_j)_{j\in J}$ de l'ensemble d'indices $I$; nous nous
proposons de donner une relation d'« associativité » entre les hypertor
$\mathfrak{Tor}_n^S((f_i)_{i\in I};(\mathscr P_\bullet^{(i)})_{i\in I})$ et
chacun des hypertor « partiels »
\[
  \mathcal T_{\bullet j}
  =\mathfrak{Tor}_\bullet^S
  ((f_i)_{i\in I_j};(\mathscr P_\bullet^{(i)})_{i\in I_j}).
\]

